\section{Proposed Method}
\label{sec:method}

\subsection{Problem and Search Strategies}
A table holds rows $(i,\mathbf{x}_i,a_i)$ with identifier $i$, embedding $\mathbf{x}_i\in\mathbb{R}^{d}$ and numeric
attribute $a_i$. A query consists of a vector $\mathbf{q}$, a range $[\ell,u]$ and $k$. With the eligible set
$E=\{i:\ell\le a_i\le u\}$ and pgvector's cosine distance
\begin{equation}
\delta(\mathbf{x},\mathbf{q})=1-\mathbf{x}^{\top}\mathbf{q}\,/\,(\lVert\mathbf{x}\rVert\,\lVert\mathbf{q}\rVert),
\label{eq:cos}
\end{equation}
the exact answer $G_k$ is the $\min(k,|E|)$ eligible rows with the smallest distance (ties broken by $i$); the
selectivity is $\sel=|E|/N$.

Every strategy issues \code{SELECT id FROM items WHERE} \emph{predicate} \code{ORDER BY embedding <=> q LIMIT 10} on
its own session with fixed settings (Table~\ref{tab:strategies}). B1 sets \code{enable\_indexscan=off}, pgvector's
documented way to obtain exact results~\cite{pgvector}: HNSW is reachable only through an index scan, so the planner
filters through the B-tree or a sequential scan and sorts every eligible row by distance. At low selectivity a B-tree
plan is cheaper, so plain SQL cannot force HNSW. B2 therefore writes the predicate as $(a+0.0)$ \code{BETWEEN} $\ell$
\code{AND} $u$, which the B-tree cannot serve, and disables sequential scans. Because $x+0.0=x$ for every finite
double, the result is unchanged. B2-iter differs from B2-plain only in iterative scanning, which continues until $k$
rows qualify or \code{max\_scan\_tuples} is reached, so it is not guaranteed to be complete. B3 runs the unmodified
query and measures PostgreSQL's own choice.

\begin{table}[t]
\caption{Search strategies and the plans verified by \code{EXPLAIN}}
\label{tab:strategies}
\centering
\footnotesize
\setlength{\tabcolsep}{3pt}
\begin{tabular}{@{}l>{\raggedright\arraybackslash}p{2.75cm}>{\raggedright\arraybackslash}p{3.65cm}@{}}
\toprule
Method & Predicate; session & Plan verified by \code{EXPLAIN}\\
\midrule
B1 & $a\in[\ell,u]$; index scans off & B-tree bitmap or (parallel) sequential scan + sort; never HNSW\\
B2-plain & $(a+0.0)\in[\ell,u]$; seq.\ scans off & HNSW index scan + filter\\
B2-iter & as B2-plain; \code{relaxed\_order} & HNSW iterative scan + filter\\
B3 & $a\in[\ell,u]$; defaults & B-tree for $\sel\le1\%$; HNSW for $\sel\ge2\%$\\
P1 & B1 or B2-iter by Eq.~\eqref{eq:route} & plan of the chosen arm\\
\bottomrule
\multicolumn{3}{@{}p{\columnwidth}@{}}{\scriptsize All HNSW sessions share \code{ef\_search=100} and
\code{max\_scan\_tuples=20000}.}
\end{tabular}
\end{table}

\subsection{Selectivity Estimation}
A \code{COUNT(*)} would cost a scan proportional to $|E|$, the very cost the router tries to avoid. P1 instead reads
PostgreSQL's \code{pg\_stats} once after \code{ANALYZE}: the null fraction $f_{\mathrm{null}}$, the most common values
$V$ with frequencies $f_v$, the equi-depth histogram bounds $b_0<\dots<b_B$, and the row count $\hat{N}$. Mirroring
PostgreSQL's range estimator~\cite{pgdocs-rowest},
\begin{equation}
\hat{P}(a\le c)=\sum_{v\in V,\,v\le c} f_v+\Bigl(1-f_{\mathrm{null}}-\sum_{v\in V} f_v\Bigr)H(c),
\label{eq:cdf}
\end{equation}
where $H(c)=\frac{1}{B}\bigl(j+\frac{c-b_j}{b_{j+1}-b_j}\bigr)$ for $b_j\le c<b_{j+1}$ interpolates inside a bucket.
The estimated number of matching rows is
\begin{equation}
\Mhat(\ell,u)=\hat{N}\,\bigl(\hat{P}(a\le u)-\hat{P}(a<\ell)\bigr).
\label{eq:est}
\end{equation}
An estimate is a few binary searches in the client, with no SQL per query.

\subsection{Routing and Threshold Calibration}
Given a frozen threshold $T$, P1 routes each query by
\begin{equation}
\pi_T(\mathbf{q},\ell,u)=
\begin{cases}
\text{B1 (exact)} & \text{if } \Mhat(\ell,u)\le T,\\
\text{B2-iter (HNSW, iterative)} & \text{otherwise.}
\end{cases}
\label{eq:route}
\end{equation}
B1's work grows with $|E|$, while B2-iter's work shrinks with $|E|$ because qualifying candidates appear early in the
traversal. A threshold on $\Mhat$ separates the two regimes if their crossover is stable across query vectors.
Fig.~\ref{fig:arch} shows the architecture; P1's latency is measured from before estimation until the last row is
fetched.

\begin{figure}[t]
\centering
\begin{tikzpicture}[
  font=\scriptsize, >={Stealth[length=1.4mm]},
  box/.style={draw, rounded corners=1.5pt, align=center, inner sep=2pt},
  off/.style={box, dashed, fill=gray!8, text width=1.55cm},
  arr/.style={->, semithick}]
\node[box, text width=3.0cm] (q)   at (0, 0)     {filtered query $(\mathbf{q},\ell,u,k{=}10)$};
\node[box, text width=3.0cm] (est) at (0,-0.85)  {estimate $\Mhat$ by Eq.~\eqref{eq:est}\\(in-process, no SQL)};
\node[draw, diamond, aspect=2.4, inner sep=0.5pt] (d) at (0,-1.85) {$\Mhat\le T^{*}$?};
\node[box, text width=2.7cm] (b1)  at (-1.95,-3.05) {B1 session (exact):\\B-tree or seq.\ scan\\+ exact distance sort};
\node[box, text width=2.7cm] (b2)  at ( 1.95,-3.05) {B2-iter session:\\HNSW iterative scan\\+ post-filter};
\node[box, text width=7.3cm] (pg)  at (0,-4.15)  {PostgreSQL 17 + pgvector 0.8.7: \code{items}$(i,a,\mathbf{x})$, B-tree on $a$, HNSW on $\mathbf{x}$};
\node[box] (out) at (0,-4.8) {top-$k$ identifiers};
\node[off] (stats) at (-3.25,-0.85) {\code{pg\_stats} snapshot after \code{ANALYZE}};
\node[off] (cal)   at ( 3.25,-1.2)  {calibration split (70\%) $\rightarrow$ frozen $T^{*}$, Eq.~\eqref{eq:calib}};
\draw[arr] (q) -- (est);
\draw[arr] (est) -- (d);
\draw[arr] (d.west) -| node[pos=0.25, above] {yes} (b1.north);
\draw[arr] (d.east) -| node[pos=0.25, above] {no} (b2.north);
\draw[arr] (b1.south) -- (b1.south |- pg.north);
\draw[arr] (b2.south) -- (b2.south |- pg.north);
\draw[arr] (pg) -- (out);
\draw[arr, dashed] (stats) -- (est);
\draw[arr, dashed] (cal.west) -- (d.north east);
\end{tikzpicture}
\caption{Architecture of the selectivity-aware router (P1). Dashed boxes are prepared offline: the statistics snapshot
is taken once after \code{ANALYZE}, and $T^{*}$ is calibrated on the calibration split and frozen before the held-out
evaluation.}
\label{fig:arch}
\end{figure}

$T$ is chosen offline. The 300 query vectors are split into 210 calibration and 90 test vectors, so all nine filters
of a vector fall on the same side. On the calibration configurations, B1 and B2-iter are measured, and P1 is simulated
for each candidate $T$ by charging every configuration the measurements of the arm Eq.~\eqref{eq:route} selects plus
the measured routing overhead. With $L(T)$, $R(T)$ and $K(T)$ the resulting mean latency, mean Recall@10 and
completeness,
\begin{equation}
T^{*}=\arg\min_{T} L(T)\quad\text{s.t.}\quad R(T)\ge0.95,\ \ K(T)\ge0.99,
\label{eq:calib}
\end{equation}
over $T=0$, $T=\infty$ and the midpoints between consecutive distinct estimates. $T^{*}$ is frozen before the held-out
run, and the HNSW settings are fixed in advance rather than tuned with $T$.
